% =========================================================================
\section{Circuito 3: amplificador inversor generalizado y filtro paso bajo}\label{sec:filtro}
% =========================================================================
\begin{prelab}[Inversor generalizado y filtro]
Sea el circuito de la \figpanel{fig:generalizado}{a}, en el que las impedancias $Z_1$ y $Z_2$ pueden ser no
resistivas y, por lo tanto, tener dependencia en frecuencia: obténgase la ganancia en voltaje del circuito.
Esta ganancia será función de la frecuencia salvo que ambas impedancias sean puramente resistivas.
Considérese el caso de que $Z_1=R$ sea real y $Z_2$ una composición en paralelo de una resistencia $R_F$ y
un condensador $C$ (\figpanel{fig:generalizado}{b}). Demuéstrese que la dependencia de la función de
transferencia con la frecuencia $f$ es la de un filtro paso bajo: para frecuencias suficientemente bajas
la ganancia en tensión es constante, mientras que para frecuencias suficientemente altas disminuye.
\end{prelab}

% Figura: (a) amplificador inversor generalizado; (b) caso Z1=R, Z2=RF||C (filtro paso bajo).
\begin{figure}[ht]
\centering
\begin{minipage}[b]{0.46\textwidth}\centering
\begin{circuitikz}[circuito, scale=0.88, transform shape]
  \draw (0,0) node[op amp, ao] (A) {};
  \draw (A.-) -- ++(-0.5,0) coordinate (N);
  \draw (N) to[generic, l_=$Z_1$, i<^=$I_1$] ++(-2.2,0) to[short, -o] ++(-0.3,0) node[left, tens] {$\Vin$};
  \draw (A.out) -- ++(0.6,0) coordinate (O);
  \draw (N) -- ++(0,1.4) coordinate (T) to[generic, l=$Z_2$, i>^=$I_2$] (T -| O) -- (O);
  \draw (O) to[short, *-o] ++(0.6,0) node[right, tens] {$\Vout$};
  \draw (A.+) -- ++(-0.5,0) -- ++(0,-0.4) node[ground] {};
  \fill (N) circle (1.6pt);
\end{circuitikz}\\[2pt]
{\small (a)}
\end{minipage}\hfill
\begin{minipage}[b]{0.50\textwidth}\centering
\begin{circuitikz}[circuito, scale=0.88, transform shape]
  \draw (0,0) node[op amp, ao] (A) {};
  \draw (A.-) -- ++(-0.5,0) coordinate (N);
  \draw (N) to[R, l_=$R$] ++(-2.2,0) to[short, -o] ++(-0.3,0) node[left, tens] {$\Vin$};
  \draw (A.out) -- ++(0.6,0) coordinate (O);
  \draw (N) -- ++(0,1.2) coordinate (T1) to[C, l_=$C$] (T1 -| O) -- (O);
  \draw (T1) -- ++(0,1.0) coordinate (T2) to[R, l=$R_F$] (T2 -| O) -- (T1 -| O);
  \fill (T1) circle (1.6pt); \fill (T1 -| O) circle (1.6pt);
  \draw (O) to[short, *-o] ++(0.6,0) node[right, tens] {$\Vout$};
  \draw (A.+) -- ++(-0.5,0) -- ++(0,-0.4) node[ground] {};
  \fill (N) circle (1.6pt);
\end{circuitikz}\\[2pt]
{\small (b)}
\end{minipage}
\caption{(a) Amplificador inversor generalizado (figura~4a del guion): $Z_1(s)$ y $Z_2(s)$
son impedancias arbitrarias. (b) Caso $Z_1=R$, $Z_2=R_F\parallel C$ (figura~5a del guion).}
\label{fig:generalizado}
\end{figure}


% -------------------------------------------------------------------------
\subsection{Herramienta: impedancia y transformada de Laplace}
% -------------------------------------------------------------------------
La ley del condensador, $i=C\,\dif v/\dif t$, es diferencial. La transformada de Laplace de una señal
$f(t)$, definida para $t\ge0$,
\begin{equation}\label{eq:L-def}
  F(s)=\Lt\{f\}(s)=\int_{0}^{\infty}f(t)\,\eul^{-st}\,\dif t ,
\end{equation}
la convierte en algebraica. Para demostrarlo, calculamos la transformada de la derivada integrando por
partes con $u=\eul^{-st}$ y $\dif w=\frac{\dif f}{\dif t}\,\dif t$, es decir, $\dif u=-s\,\eul^{-st}\dif t$ y $w=f(t)$:
\begin{equation}\label{eq:L-partes}
  \Lt\Bigl\{\frac{\dif f}{\dif t}\Bigr\}=\int_{0}^{\infty}\frac{\dif f}{\dif t}\,\eul^{-st}\,\dif t
  =\Bigl[f(t)\,\eul^{-st}\Bigr]_{0}^{\infty}-\int_{0}^{\infty}f(t)\,\bigl(-s\,\eul^{-st}\bigr)\,\dif t .
\end{equation}
En el límite superior, $f(t)\,\eul^{-st}\to0$ cuando $t\to\infty$ (las señales del circuito están acotadas y
$\Real s>0$); en el inferior, $\eul^{0}=1$. El término de contorno vale, por tanto, $0-f(0)$, y la integral
restante es $s$ veces \eqref{eq:L-def}:
\begin{equation}\label{eq:L-deriv0}
  \Lt\Bigl\{\frac{\dif f}{\dif t}\Bigr\}=s\int_{0}^{\infty}f(t)\,\eul^{-st}\,\dif t-f(0)=sF(s)-f(0) .
\end{equation}
Con condición inicial nula, $f(0)=0$ (condensador descargado), derivar equivale a multiplicar por $s$.
Aplicando la transformada, que es lineal, a los dos miembros de $i=C\,\dif v/\dif t$ con $f=v$:
\begin{equation}\label{eq:L-deriv}
  I(s)=C\,\Lt\Bigl\{\frac{\dif v}{\dif t}\Bigr\}=C\,\bigl[sV(s)-v(0)\bigr]
  \overset{v(0)=0}{=}sC\,V(s) .
\end{equation}
Cada elemento cumple así una «ley de Ohm» $V(s)=Z(s)\,I(s)$, con impedancias
\begin{equation}\label{eq:L-Z}
  Z_R=R,\qquad Z_C(s)=\frac{1}{sC},
\end{equation}
que se asocian como resistencias. Por linealidad, Kirchhoff, (R1) y (R2) valen para las transformadas, y en
régimen sinusoidal de pulsación $\omega=2\pi f$ basta tomar $s=\jj\omega$: $\abs{H(\jj\omega)}$ es el cociente de
amplitudes $\Vout/\Vin$ y $\arg H(\jj\omega)$ el desfase.

% -------------------------------------------------------------------------
\subsection{Ganancia del inversor generalizado}
% -------------------------------------------------------------------------
\begin{resolucion}
Se repite el cálculo del inversor con $R_k\to Z_k(s)$ (sentidos de la \figpanel{fig:generalizado}{a}).

\textbf{1. Tierra virtual} (R2): la entrada $+$ está a tierra,
\begin{equation}\label{eq:g-vm}
  V_-(s)=V_+(s)=0 .
\end{equation}

\textbf{2. Ohm en $Z_1$ y Kirchhoff en el nudo inversor}, con $I_-=0$ (R1):
\begin{equation}\label{eq:g-I}
  I_2(s)=I_1(s)=\frac{\Vin(s)-V_-(s)}{Z_1(s)}=\frac{\Vin(s)}{Z_1(s)} .
\end{equation}

\textbf{3. Ohm en $Z_2$}, del nudo inversor a la salida:
\begin{equation}\label{eq:g-vout}
  V_-(s)-\Vout(s)=Z_2(s)\,I_2(s)\quad\Longrightarrow\quad\Vout(s)=-\frac{Z_2(s)}{Z_1(s)}\,\Vin(s) .
\end{equation}

\textbf{Resultado.}
\begin{equation}\label{eq:g-H}
  \resultado{H(s)=\frac{\Vout(s)}{\Vin(s)}=-\frac{Z_2(s)}{Z_1(s)}},\qquad
  A_V(\jj\omega)=-\frac{Z_2(\jj\omega)}{Z_1(\jj\omega)} .
\end{equation}
Si $Z_1=R_1$ y $Z_2=R_2$, $A_V=-R_2/R_1$ es real y no depende de $\omega$; si hay un condensador,
$Z_C=1/(\jj\omega C)$ introduce la dependencia en frecuencia.
\end{resolucion}

% -------------------------------------------------------------------------
\subsection{Caso $Z_1=R$, $Z_2=R_F\parallel C$: filtro paso bajo}
% -------------------------------------------------------------------------
\begin{resolucion}
\textbf{1. Impedancias.} Por \eqref{eq:L-Z}, sumando las inversas del paralelo,
\begin{equation}\label{eq:f-Z2}
  \frac{1}{Z_2}=\frac{1}{R_F}+sC=\frac{1+sR_FC}{R_F}\quad\Longrightarrow\quad
  Z_1=R,\qquad Z_2(s)=\frac{R_F}{1+sR_FC} .
\end{equation}

\textbf{2. Función de transferencia.} Sustituyendo \eqref{eq:f-Z2} en \eqref{eq:g-H}:
\begin{equation}\label{eq:f-H}
  \resultado{H(s)=-\frac{R_F}{R}\,\frac{1}{1+sR_FC}},\qquad
  f_c=\frac{1}{2\pi R_FC} .
\end{equation}

\textbf{3. Módulo de la ganancia.} Con $s=\jj\omega=\jj2\pi f$ y $\abs{1+\jj x}=\sqrt{1+x^{2}}$,
\begin{equation}\label{eq:f-mod}
  \abs{A_V(f)}=\frac{R_F/R}{\sqrt{1+(f/f_c)^{2}}},\qquad
  \arg A_V(f)=180^{\circ}-\arctg\frac{f}{f_c} .
\end{equation}

\textbf{4. Límites.} De \eqref{eq:f-mod}:
\begin{align}
  f\ll f_c:&\quad \abs{A_V}\simeq\frac{R_F}{R}\quad\text{(constante; $C$ actúa como circuito abierto)},\label{eq:f-baja}\\
  f\gg f_c:&\quad \abs{A_V}\simeq\frac{R_F}{R}\,\frac{f_c}{f}=\frac{1}{2\pi fRC}\quad\text{(decrece como $1/f$; $C$ cortocircuita a $R_F$)}.\label{eq:f-alta}
\end{align}
En decibelios, por \eqref{eq:f-alta}, cada década de frecuencia resta
$20\log_{10}10=20\dB$; y en $f=f_c$, $\abs{A_V}=(R_F/R)/\sqrt2$, es decir, $3\dB$ por debajo de la meseta.

\textbf{Conclusión.} La ganancia es constante a baja frecuencia \eqref{eq:f-baja} y disminuye a
$-20\dB$/década a alta frecuencia \eqref{eq:f-alta}, con la transición en $f_c$: el circuito es un
\emph{filtro paso bajo de primer orden}, como en la figura~4b del guion (\cref{fig:bode-filtro}).
\end{resolucion}

% Figura: Bode del filtro paso bajo activo, con las frecuencias características en forma simbólica.
\begin{figure}[ht]
\centering
\begin{tikzpicture}
  \begin{semilogxaxis}[tema, name=mod, width=0.48\textwidth, height=5.4cm,
      xmin=10, xmax=1e6, ymin=-12, ymax=48, xtick={1001,100097}, xticklabels={$f_c$,$f_1$},
      ytick={0,37,40}, yticklabels={$0$,,$20\log\frac{R_F}{R}$},
      xlabel={$f$ (escala logarítmica)}, ylabel={$\abs{A_V}$ (dB)}, log basis x=10,
      title style={font=\small}, title={(a) Módulo}, clip=false]
    \addplot[asintota, domain=10:1001, samples=2] {40};
    \addplot[asintota, domain=1001:100097, samples=2] {20*log10(100097/x)};
    \addplot[azulacero, line width=1.1pt, domain=10:1e6, samples=300]
      {40-10*log10(1+(x/1001)^2)};
    \addplot[only marks, mark=*, mark size=1.6pt, rojoteo] coordinates {(1001,37)};
    \node[font=\scriptsize, anchor=north east, text=rojoteo] at (axis cs:950,35) {$-3\dB$};
    \node[font=\scriptsize, anchor=east, text=grispizarra] at (axis cs:2500,12) {$-20\dB$/década $\rightarrow$};
  \end{semilogxaxis}
  \begin{semilogxaxis}[tema, at={(mod.east)}, anchor=west, xshift=1.6cm,
      width=0.48\textwidth, height=5.4cm,
      xmin=10, xmax=1e6, ymin=80, ymax=190, ytick={90,135,180},
      xtick={1001}, xticklabels={$f_c$},
      xlabel={$f$ (escala logarítmica)}, ylabel={$\arg A_V$ (grados)}, log basis x=10,
      title style={font=\small}, title={(b) Fase}]
    \addplot[azulacero, line width=1.1pt, domain=10:1e6, samples=300] {180-atan(x/1001)};
    \addplot[only marks, mark=*, mark size=1.6pt, rojoteo] coordinates {(1001,135)};
  \end{semilogxaxis}
\end{tikzpicture}
\caption{Diagrama de Bode del filtro de la \figpanel{fig:generalizado}{b}. (a) Módulo, según
\eqref{eq:f-dB}: meseta $20\log_{10}(R_F/R)$, corte a $-3\dB$ en $f_c=1/(2\pi R_FC)$ y caída de
$-20\dB$/década (asíntota discontinua, que corta $0\dB$ en $f_1=1/(2\pi RC)$). (b) Fase, según
\eqref{eq:f-mod}: de $180^{\circ}$ a $90^{\circ}$, con $135^{\circ}$ en $f_c$.}
\label{fig:bode-filtro}
\end{figure}

